% Card: the long side (rook endings)
\paragraph{Long side}\index{Rook endings!Long side}

Every pawn except a rook's pawn divides the board into two parts of different width. The wider
part is the
\vocab[rook endings]{long-side}{long side}{The wider of the two parts into
which a pawn (other than a rook's pawn) divides the board. The defending rook belongs there,
so that it has room for side checks}.
In \xref{diagram}{side-checks} the e-pawn has four files on its left (a, b, c, d): that is
the long side.

Because \vocab{side-checks}{side checks}{} need room, the defending rook belongs on the long
side. Put Black's rook on h2 instead of a2 in \xref{diagram}{side-checks}: it is too close to
the pawn to give side checks, and White wins.
